\chapter{Galerkin's Method and the Bilinear-Form Family}
\label{ch:galerkin}

\chapterorient{In \cref{ch:weak} we turned a partial differential equation into
a \emph{weak form}: find a function $u$ such that $a(u,v)=\ell(v)$ for every test
function $v$. That statement is exact, but it lives in an infinite-dimensional
space of functions---there is no way to hand it to a computer as written. This
chapter performs the single decisive step that makes the weak form
\emph{computable}: replace the infinite-dimensional space with a finite one
spanned by $N$ chosen functions, and the weak form collapses into a system of
$N$ linear equations in $N$ unknowns, $\mat{K}\vx=\vb$. We will see where the
matrices $\mat{K}$ and $\mat{M}$ that the rest of the book solves actually
\emph{come from}---they are nothing but the bilinear form evaluated on pairs of
basis functions---and why the properties an engineer cares about (symmetry,
sparsity, positive-definiteness) are inherited directly from the form. By the
end you will have assembled, by hand, the classic tridiagonal stiffness matrix
of the one-dimensional Poisson equation.}

\section{From an infinite problem to a finite one}

Recall the shape of the weak form. We have a real number $a(u,v)$ that depends
linearly on each of two functions $u$ and $v$---a \emph{bilinear form}---and a
real number $\ell(v)$ that depends linearly on one function---a \emph{linear
functional}. The weak problem is:

\begin{quote}
find $u$ in the trial space $V$ such that $a(u,v)=\ell(v)$ for \emph{all} $v$ in
the test space $V$.
\end{quote}

For the model problem of this chapter, the one-dimensional Poisson equation
$-u''=f$ on the interval $(0,1)$ with $u(0)=u(1)=0$, integration by parts gave us
(in \cref{ch:weak}) the concrete pair
\begin{equation}
a(u,v)=\int_0^1 u'(x)\,v'(x)\dx,
\qquad
\ell(v)=\int_0^1 f(x)\,v(x)\dx.
\label{eq:poisson-weak}
\end{equation}
The space $V$ here is the set of all functions on $(0,1)$ that vanish at the two
endpoints and whose first derivative is square-integrable---an
\emph{infinite-dimensional} space. We cannot loop over ``all $v$ in $V$''; there
are uncountably many, and no finite list of equations pins $u$ down. Something
has to give.

\subsection{The Galerkin idea}

The idea, due to Boris Galerkin, is disarmingly simple: \emph{stop asking for the
exact answer in the infinite space, and instead ask for the best answer you can
build out of a finite menu of functions.} We pick, once and for all, $N$
functions
\[
\phi_1,\phi_2,\dots,\phi_N,
\]
each one a legal member of $V$ (each vanishes at the boundary, each is smooth
enough), and we agree to look for our solution only among combinations of them.
The set of all such combinations is a finite-dimensional subspace,
\begin{equation}
V_h=\operatorname{span}\{\phi_1,\dots,\phi_N\}
=\Bigl\{\,\textstyle\sum_{j=1}^N c_j\phi_j : c_j\in\Real\,\Bigr\}\subset V,
\label{eq:Vh}
\end{equation}
where the subscript $h$ is the traditional reminder that the menu came from a
\emph{mesh} of spacing $h$ (we make this concrete in \cref{sec:worked-stiffness}).
The functions $\phi_i$ are the \emph{basis functions}, and choosing them well is,
as \cref{sec:basis-is-everything} argues, the entire game.

\begin{figure}[t]
\centering
\begin{tikzpicture}[scale=1.0, every node/.style={font=\footnotesize}]
  % the big infinite space V
  \draw[draw=simGray, fill=simBgGray, rounded corners=10pt]
    (-3.4,-1.9) rectangle (3.4,2.0);
  \node[simGray] at (0,1.65) {$V$ \;(infinite-dimensional)};
  % the finite subspace V_h
  \draw[draw=simGreen, fill=simBgGreen, rounded corners=6pt, thick]
    (-2.6,-1.5) rectangle (2.6,0.7);
  \node[simGreen!70!simInk] at (0,0.42) {$V_h=\operatorname{span}\{\phi_1,\dots,\phi_N\}$};
  % the exact solution
  \fill[simRust] (1.9,1.25) circle (1.6pt);
  \node[simRust, above right=-1pt and -2pt] at (1.9,1.25) {$u$ (exact)};
  % the galerkin solution inside V_h
  \fill[simBlue] (0.55,-0.55) circle (1.6pt);
  \node[simBlue, below=1pt] at (0.55,-0.7) {$u_h$ (Galerkin)};
  % projection arrow
  \draw[flow, simInk] (1.9,1.25) -- node[right=1pt,font=\scriptsize]{project} (0.55,-0.55);
  % the residual is orthogonal to V_h: little right-angle marker idea
  \node[font=\scriptsize\itshape, simInk] at (-1.45,-1.05)
    {$a(u-u_h,\,v_h)=0\ \ \forall v_h\in V_h$};
\end{tikzpicture}
\caption[The Galerkin projection]{The Galerkin idea in one picture. The exact
solution $u$ lives somewhere in the infinite-dimensional space $V$; we cannot
reach it. Instead we look inside the finite-dimensional subspace $V_h$ and pick
the element $u_h$ whose \emph{error} $u-u_h$ is invisible to every test function
in $V_h$---formally, $a(u-u_h,v_h)=0$ for all $v_h\in V_h$. In the energy
geometry set by $a(\cdot,\cdot)$ this is an orthogonal projection, which is why
$u_h$ turns out to be the \emph{best} approximation $V_h$ can offer
(\cref{sec:cea}).}
\label{fig:galerkin-projection}
\end{figure}

Having narrowed the search to $V_h$, we must also narrow the supply of test
functions: it is unreasonable to demand $a(u_h,v)=\ell(v)$ for \emph{all} $v$ in
the infinite $V$ when $u_h$ only has $N$ knobs to turn. The \emph{Galerkin
choice} is to use the same finite menu for both roles---trial and test---and
require the weak equation only against the $N$ basis functions:
\begin{equation}
\text{find } u_h\in V_h \text{ such that }
a(u_h,\phi_i)=\ell(\phi_i)
\quad\text{for } i=1,2,\dots,N.
\label{eq:galerkin}
\end{equation}
Because $a$ and $\ell$ are linear in their test slot, requiring
\cref{eq:galerkin} against each basis function $\phi_i$ is exactly the same as
requiring it against every combination $v_h=\sum_i d_i\phi_i$---so
\cref{eq:galerkin} really does enforce the weak form for \emph{all} of $V_h$,
not just the basis. We have replaced ``for all $v\in V$,'' an infinitude of
conditions, with ``for $i=1,\dots,N$,'' exactly $N$ of them.

\subsection{Counting: \texorpdfstring{$N$}{N} equations, \texorpdfstring{$N$}{N} unknowns}

Now write the unknown $u_h$ out in the basis. Every element of $V_h$ is some
combination $\sum_j u_j\phi_j$; finding $u_h$ \emph{means} finding the
coefficients $u_1,\dots,u_N$. So we substitute
\begin{equation}
u_h=\sum_{j=1}^N u_j\,\phi_j
\label{eq:uh-expansion}
\end{equation}
into the Galerkin condition \cref{eq:galerkin} and use bilinearity of $a$ to pull
the unknown coefficients out of the form:
\begin{equation}
a(u_h,\phi_i)
=a\Bigl(\sum_{j=1}^N u_j\phi_j,\;\phi_i\Bigr)
=\sum_{j=1}^N u_j\,a(\phi_j,\phi_i)
=\ell(\phi_i),
\qquad i=1,\dots,N.
\label{eq:galerkin-expanded}
\end{equation}
Look at what \cref{eq:galerkin-expanded} is. For each $i$ it is one linear
equation in the unknowns $u_1,\dots,u_N$. There are $N$ values of $i$, hence $N$
equations; there are $N$ unknowns $u_j$. The numbers $a(\phi_j,\phi_i)$ and
$\ell(\phi_i)$ are not unknown---they are definite integrals we can compute from
the data of the problem. The infinite-dimensional weak form has become a
square linear system. That is the whole of Galerkin's method.

\begin{keyidea}
The matrices \emph{are} the bilinear form, sampled on basis pairs. There is no
separate ``matrix-building'' physics: the entry in row $i$, column $j$ of the
stiffness matrix is \emph{defined} to be $a(\phi_j,\phi_i)$, the bilinear form
fed the $j$th and $i$th basis functions. Symmetry, definiteness, sparsity---every
property of the matrix is a property of the form $a$ read off through this
sampling. To know the matrix, know the form and know the basis; nothing else
enters.
\end{keyidea}

\section{The assembled system \texorpdfstring{$\mat{K}\vx=\vb$}{Kx=b}}

We collect \cref{eq:galerkin-expanded} into matrix--vector form. Define the
$N\times N$ matrix $\mat{K}$ and the $N$-vectors $\vx,\vb$ by
\begin{equation}
\mat{K}_{ij}=a(\phi_j,\phi_i),
\qquad
\vx_j=u_j,
\qquad
\vb_i=\ell(\phi_i).
\label{eq:stiffness-def}
\end{equation}
Then the $N$ Galerkin equations \cref{eq:galerkin-expanded} are precisely the
single matrix equation
\begin{equation}
\boxed{\;\mat{K}\,\vx=\vb\;}
\label{eq:Kxb}
\end{equation}
This is the discrete problem the rest of the book solves. The matrix $\mat{K}$ is
called the \emph{stiffness matrix}, a name inherited from structural mechanics,
where $a(u,v)$ measured elastic energy and its discretization measured how stiffly
a structure resisted deformation. The vector $\vb$ is the \emph{load vector}, and
$\vx$ holds the coefficients of the discrete solution $u_h=\sum_j x_j\phi_j$.

\begin{notationbox}
A small but constant source of confusion: the index order in
$\mat{K}_{ij}=a(\phi_j,\phi_i)$ has $j$ (the trial index, the column) inside the
form's \emph{first} slot and $i$ (the test index, the row) in the \emph{second}.
We chose this so that $\mat{K}\vx$ produces row $i=\sum_j a(\phi_j,\phi_i)x_j
=a(u_h,\phi_i)$, the left side of the Galerkin equation. When $a$ is symmetric
($a(u,v)=a(v,u)$, the case for all the stiffness and mass matrices in this book)
the order does not matter and $\mat{K}_{ij}=\mat{K}_{ji}$. It matters only for the
\emph{non}symmetric forms we meet in the driven analysis (\cref{ch:driven}); we
flag it now so the convention is never a surprise.
\end{notationbox}

\subsection{Where \texorpdfstring{$\mat{K}$}{K} and \texorpdfstring{$\mat{M}$}{M} come from}

The two matrices that dominate this book---the stiffness $\mat{K}$ and the
\emph{mass} $\mat{M}$---are both instances of \cref{eq:stiffness-def}, differing
only in which bilinear form is sampled. The stiffness matrix samples a form built
from \emph{derivatives} of the basis functions; the mass matrix samples the
plain $L^2$ pairing with \emph{no} derivative:
\begin{equation}
\mat{K}_{ij}=a(\phi_j,\phi_i)=\int_\dom Q\,\diffop\phi_j\cdot\diffop\phi_i\dV,
\qquad
\mat{M}_{ij}=\int_\dom \phi_i\,\phi_j\dV.
\label{eq:KM}
\end{equation}
Here $\diffop$ is whichever differential operator the weak form attached to the
trial and test functions---$\Grad$, $\Curl$, $\Div$, or the identity---and $Q$ is
a material weight (a permittivity $\varepsilon$, a reluctivity $\nu=1/\mu$, and so
on). The mass matrix is the special case $\diffop=\text{identity}$ and $Q=1$:
$\mat{M}_{ij}=\int\phi_i\phi_j$ measures the overlap of basis function $i$ with
basis function $j$.

We will meet these two matrices in every analysis. The eigenmode problem of
\cref{ch:eigenmodes} is the generalized eigenvalue problem $\mat{K}\vx=\lambda
\mat{M}\vx$; the transient problem of \cref{ch:transient} marches
$\mat{M}\ddot{\vs}+\mat{C}\dot{\vs}+\mat{K}\vs=\vJ(t)$; the driven problem of
\cref{ch:driven} solves $(\mat{K}+i\omega\mat{C}-\omega^2\mat{M})\vE=\vb$. In
\emph{every} one of these, $\mat{K}$ and $\mat{M}$ are the bilinear forms of
\cref{eq:KM} sampled on basis pairs. They look like distinct objects in a linear
algebra library; they are, in truth, the same form-on-basis construction with
different $\diffop$ and $Q$.

\begin{physicsbox}
The mass matrix earns its name from dynamics. In the wave-equation
discretization $\mat{M}\ddot{\vs}+\mat{K}\vs=0$, the term $\mat{M}\ddot{\vs}$ is
the discrete analogue of ``mass times acceleration.'' Off-diagonal entries
$\mat{M}_{ij}\neq0$ mean basis functions $i$ and $j$ overlap, so that
accelerating the coefficient $u_j$ exerts an inertial pull on equation $i$---the
finite-element mesh has, in effect, a smeared-out mass distribution rather than
point masses. A diagonal (``lumped'') mass matrix is the point-mass
idealization; the true \emph{consistent} mass matrix of \cref{eq:KM} couples
neighbours, as we will see in \cref{wex:mass}.
\end{physicsbox}

\section{Three properties that decide everything downstream}

The matrix $\mat{K}$ inherits its character from the form $a$. Three inherited
properties govern which solver the rest of the book reaches for, and we take them
in turn.

\subsection{Symmetry}

If the form is symmetric, $a(u,v)=a(v,u)$, then by \cref{eq:stiffness-def}
\[
\mat{K}_{ij}=a(\phi_j,\phi_i)=a(\phi_i,\phi_j)=\mat{K}_{ji},
\]
so $\mat{K}$ is a symmetric matrix, $\mat{K}=\mat{K}^{\mathsf T}$. Every form of
the family \cref{eq:KM} with a symmetric (in particular, scalar) weight $Q$ is
symmetric, because $Q\,\diffop\phi_j\cdot\diffop\phi_i=Q\,\diffop\phi_i\cdot
\diffop\phi_j$ pointwise. Symmetry is not a cosmetic nicety: it is the hypothesis
the conjugate-gradient method requires. When $\mat{K}$ is symmetric (and positive
definite, below), the system $\mat{K}\vx=\vb$ can be solved by conjugate
gradients, the workhorse iterative method of \cref{ch:cg}---which exploits exactly
this symmetry to build its short three-term recurrence. A nonsymmetric $\mat{K}$
forces the more expensive GMRES of \cref{ch:gmres} instead. The choice of solver
is decided here, at assembly time, by a property of the continuous form.

\subsection{Coercivity and positive-definiteness}

\Cref{ch:weak} introduced \emph{coercivity}: there is a constant $\alpha>0$ with
$a(v,v)\ge\alpha\norm{v}^2$ for every $v\in V$. Coercivity is what guaranteed the
continuous weak problem had a unique solution (the Lax--Milgram theorem). It
descends to the matrix as \emph{positive-definiteness}. Take any nonzero
coefficient vector $\vy=(y_1,\dots,y_N)$ and let $v_h=\sum_j y_j\phi_j$, a nonzero
element of $V_h$. Then
\begin{equation}
\vy^{\mathsf T}\mat{K}\vy
=\sum_{i,j}y_i\,\mat{K}_{ij}\,y_j
=\sum_{i,j}y_i\,a(\phi_j,\phi_i)\,y_j
=a\Bigl(\textstyle\sum_j y_j\phi_j,\sum_i y_i\phi_i\Bigr)
=a(v_h,v_h)
\ge\alpha\norm{v_h}^2>0.
\label{eq:spd}
\end{equation}
The chain \cref{eq:spd} is worth pausing on: the algebraic quantity
$\vy^{\mathsf T}\mat{K}\vy$, the thing a numerical analyst calls the matrix's
quadratic form, is \emph{equal} to the continuous energy $a(v_h,v_h)$ of the
function $v_h$ those coefficients represent. Coercivity of the form is thus,
literally, positive-definiteness of the matrix. A symmetric positive-definite
(SPD) matrix is invertible---so the discrete system has a unique solution---and it
is precisely the class on which conjugate gradients is guaranteed to converge.
Two properties an engineer wants, solvability and a fast solver, both flow from
the single analytic fact that $a$ is coercive.

\subsection{Sparsity}
\label{sec:sparsity}

The third inherited property is the one that makes large problems tractable at
all. The entry $\mat{K}_{ij}=\int_\dom Q\,\diffop\phi_j\cdot\diffop\phi_i\dV$ is a
product of two basis functions integrated over the domain. If $\phi_i$ and
$\phi_j$ are nonzero on \emph{disjoint} regions of the domain, their product is
zero everywhere and so is the integral: $\mat{K}_{ij}=0$. Finite-element basis
functions are built to have \emph{local support}---each $\phi_i$ is nonzero only
on the handful of mesh cells touching node $i$ (\cref{fig:hats})---so for the
overwhelming majority of index pairs $(i,j)$ the supports miss each other and the
entry vanishes.

\begin{figure}[t]
\centering
\begin{tikzpicture}[scale=1.0, every node/.style={font=\footnotesize}]
  % nodes around node i, showing only i-1, i, i+1 couple
  \def\R{0.16}
  \foreach \k/\x in {{i-2}/0,{i-1}/1.6,{i}/3.2,{i+1}/4.8,{i+2}/6.4}{
    \fill[simInk] (\x,0) circle (\R);
    \node[below=2pt] at (\x,0) {$\phi_{\k}$};
  }
  % highlight phi_i support region
  \draw[draw=simBlue, fill=simBgBlue, rounded corners=4pt, opacity=0.6]
    (1.6,-0.5) rectangle (4.8,0.95);
  \node[simBlue!70!simInk, above] at (3.2,0.62) {support of $\phi_i$};
  % coupling arcs: i couples only to i-1 and i+1
  \draw[depedge, bend left=40] (3.2,0.12) to (1.6,0.12);
  \draw[depedge, bend left=40] (3.2,0.12) to (4.8,0.12);
  \node[simBlue, font=\scriptsize] at (2.4,0.95) {couples};
  % the far ones don't couple: dotted cross
  \draw[refedge] (3.2,-0.28) -- (6.4,-0.28);
  \node[simGray, font=\scriptsize] at (5.4,-0.5) {$\mat{K}_{i,\,i+2}=0$};
\end{tikzpicture}
\caption[Local support gives sparsity]{Why $\mat{K}$ is sparse. A finite-element
basis function $\phi_i$ is nonzero only on the cells immediately around node $i$
(shaded). The entry $\mat{K}_{ij}=\int Q\,\diffop\phi_j\cdot\diffop\phi_i$ is
nonzero only when the supports of $\phi_i$ and $\phi_j$ \emph{overlap}, which
happens only for the few $j$ that are mesh-neighbours of $i$. Here $\phi_i$
couples to $\phi_{i-1}$ and $\phi_{i+1}$ but not to $\phi_{i+2}$, whose support is
disjoint---so that matrix entry is exactly zero.}
\label{fig:sparsity-support}
\end{figure}

The consequence is that $\mat{K}$ has only $\mathcal{O}(N)$ nonzero entries, not
$N^2$. For a one-dimensional mesh with the hat functions of
\cref{sec:worked-stiffness} each row has at most three nonzeros and $\mat{K}$ is
\emph{tridiagonal}; in three dimensions each row has a few tens of nonzeros
regardless of how large $N$ grows. A matrix with $\mathcal{O}(N)$ nonzeros can be
\emph{multiplied} by a vector in $\mathcal{O}(N)$ work without ever forming the
$N^2$ dense array---and that is the only operation the iterative solvers of
\cref{ch:krylov} ever ask of it. Sparsity is the reason we solve $\mat{K}\vx=\vb$
by iteration rather than by Gaussian elimination: the matrix is too large to store
densely but cheap to apply.

\begin{pitfall}
Sparsity is a property of the \emph{basis}, not of the differential equation.
Choose global basis functions---say $\phi_k(x)=\sin(k\pi x)$, each nonzero across
the whole interval---and \emph{every} pair of supports overlaps, so $\mat{K}$ is
dense even for the same Poisson problem. (For that particular global basis
$\mat{K}$ happens to come out diagonal by orthogonality, but that is a lucky
accident of the sine basis, not a general phenomenon.) The finite-element method
earns its sparsity precisely by insisting on \emph{locally} supported basis
functions; a different choice forfeits it. This is the first hint of the theme of
\cref{sec:basis-is-everything}: the basis is the whole game.
\end{pitfall}

\section{Worked example: the 1-D Poisson stiffness matrix}
\label{sec:worked-stiffness}

We now assemble $\mat{K}$ explicitly for the model problem
\cref{eq:poisson-weak}, and discover the most famous matrix in numerical analysis.

\subsection{The mesh and the hat basis}

Partition the interval $[0,1]$ into $n$ equal cells with $n+1$ nodes
\[
x_0=0,\quad x_1=h,\quad x_2=2h,\quad\dots,\quad x_n=1,
\qquad h=\frac1n.
\]
At each \emph{interior} node $x_i$ ($i=1,\dots,n-1$) place the \emph{hat
function} $\phi_i$: the piecewise-linear function equal to $1$ at $x_i$, equal to
$0$ at every other node, and linear in between (\cref{fig:hats}). Explicitly,
\begin{equation}
\phi_i(x)=
\begin{cases}
\dfrac{x-x_{i-1}}{h}, & x_{i-1}\le x\le x_i,\\[2ex]
\dfrac{x_{i+1}-x}{h}, & x_i\le x\le x_{i+1},\\[1ex]
0, & \text{otherwise.}
\end{cases}
\label{eq:hat}
\end{equation}
There are $N=n-1$ of them, one per interior node; the two boundary nodes carry no
unknown because $u(0)=u(1)=0$ is built in (those are the \emph{essential} boundary
conditions of \cref{ch:weak}, and we simply omit their hats). Each hat is nonzero
on exactly two cells, so its support is local and \cref{sec:sparsity} promises a
sparse---in fact tridiagonal---matrix.

\begin{figure}[t]
\centering
\begin{tikzpicture}[scale=1.0]
\begin{axis}[
  width=12cm, height=5.2cm,
  xmin=-0.05, xmax=1.05, ymin=-0.08, ymax=1.18,
  xtick={0,0.2,0.4,0.6,0.8,1.0},
  xticklabels={$x_0$,$x_1$,$x_2$,$x_3$,$x_4$,$x_5$},
  ytick={0,1}, ylabel={}, axis lines=left,
  tick label style={font=\footnotesize},
  clip=false,
]
  % hat phi_1 at x=0.2
  \addplot[simBlue, very thick] coordinates {(0,0)(0.2,1)(0.4,0)(1.0,0)};
  % hat phi_2 at x=0.4
  \addplot[simRust, very thick] coordinates {(0,0)(0.2,0)(0.4,1)(0.6,0)(1.0,0)};
  % hat phi_3 at x=0.6
  \addplot[simGreen, very thick] coordinates {(0,0)(0.4,0)(0.6,1)(0.8,0)(1.0,0)};
  % hat phi_4 at x=0.8
  \addplot[simViolet, very thick] coordinates {(0,0)(0.6,0)(0.8,1)(1.0,0)};
  % node markers
  \addplot[only marks, mark=*, mark size=1.4pt, simInk]
    coordinates {(0,0)(0.2,0)(0.4,0)(0.6,0)(0.8,0)(1.0,0)};
  \node[font=\scriptsize, simBlue]   at (axis cs:0.2,1.08) {$\phi_1$};
  \node[font=\scriptsize, simRust]   at (axis cs:0.4,1.08) {$\phi_2$};
  \node[font=\scriptsize, simGreen]  at (axis cs:0.6,1.08) {$\phi_3$};
  \node[font=\scriptsize, simViolet] at (axis cs:0.8,1.08) {$\phi_4$};
\end{axis}
\end{tikzpicture}
\caption[The hat basis]{The hat (piecewise-linear ``tent'') basis on a uniform
mesh of $n=5$ cells, $h=\tfrac15$. Each interior node $x_i$ carries one basis
function $\phi_i$ that peaks at $1$ on its own node and falls linearly to $0$ at
the two neighbouring nodes, vanishing everywhere else. The two boundary nodes
$x_0,x_5$ carry no function because $u$ is pinned to zero there. Adjacent hats
overlap on one cell (so they couple); hats two nodes apart do not (so their matrix
entry is zero)---the picture of sparsity from \cref{fig:sparsity-support}.}
\label{fig:hats}
\end{figure}

\subsection{Computing the entries}

The form is $a(\phi_j,\phi_i)=\int_0^1\phi_j'(x)\,\phi_i'(x)\dx$. The derivative
of a hat is a pair of constants: from \cref{eq:hat},
\[
\phi_i'(x)=
\begin{cases}
+1/h, & x_{i-1}<x<x_i,\\
-1/h, & x_i<x<x_{i+1},\\
0, & \text{otherwise,}
\end{cases}
\]
a $+1/h$ slope going up and a $-1/h$ slope coming down. We need three kinds of
entry.

\medskip
\noindent\textbf{Diagonal, $\mathbf{K_{ii}}$.} On its support, $\phi_i$ has slope
$\pm1/h$ over two cells each of length $h$:
\[
\mat{K}_{ii}=\int_{x_{i-1}}^{x_{i+1}}\bigl(\phi_i'\bigr)^2\dx
=\Bigl(\tfrac1h\Bigr)^2\!\cdot h+\Bigl(-\tfrac1h\Bigr)^2\!\cdot h
=\frac1h+\frac1h=\frac2h.
\]

\noindent\textbf{First off-diagonal, $\mathbf{K_{i,i+1}}$.} The hats $\phi_i$ and
$\phi_{i+1}$ overlap only on the single cell $[x_i,x_{i+1}]$. There
$\phi_i'=-1/h$ (coming down) while $\phi_{i+1}'=+1/h$ (going up), so
\[
\mat{K}_{i,i+1}=\int_{x_i}^{x_{i+1}}\phi_i'\,\phi_{i+1}'\dx
=\Bigl(-\tfrac1h\Bigr)\Bigl(+\tfrac1h\Bigr)\cdot h=-\frac1h.
\]
By symmetry $\mat{K}_{i+1,i}=-1/h$ as well.

\medskip
\noindent\textbf{Everything else, $\mathbf{K_{ij}}$ with $|i-j|\ge2$.} Disjoint
supports, so the integral is zero.

\medskip
Assembling these into the $N\times N$ matrix gives the celebrated tridiagonal
form
\begin{equation}
\mat{K}=\frac1h
\begin{pmatrix}
2 & -1 & & & \\
-1 & 2 & -1 & & \\
 & -1 & 2 & \ddots & \\
 & & \ddots & \ddots & -1\\
 & & & -1 & 2
\end{pmatrix}
=\frac1h\,\operatorname{tridiag}(-1,\,2,\,-1).
\label{eq:tridiag}
\end{equation}
This is the discrete second derivative, the same matrix a finite-difference
scheme produces by approximating $-u''$ with
$\bigl(-u_{i-1}+2u_i-u_{i+1}\bigr)/h^2$ (the $1/h$ versus $1/h^2$ difference is
just where the cell width is booked; the finite-difference RHS carries an extra
$h$). The finite-element method has reproduced the classical stencil---but it
arrived by sampling a bilinear form on basis functions, which is what lets the
same machinery handle curved elements, anisotropic materials, and vector fields
that the finite-difference stencil cannot.

\begin{figure}[t]
\centering
\begin{tikzpicture}[scale=0.62, every node/.style={font=\scriptsize}]
  \def\nz{simBlue}
  % draw a 7x7 grid; nonzeros on tridiagonal
  \foreach \r in {0,...,6}{
    \foreach \c in {0,...,6}{
      \pgfmathtruncatemacro{\d}{abs(\r-\c)}
      \ifnum\d<2
        \fill[\nz] (\c,-\r) rectangle ++(0.92,-0.92);
      \else
        \draw[simGray!40] (\c,-\r) rectangle ++(0.92,-0.92);
      \fi
    }
  }
  % frame
  \draw[simInk, thick] (0,0) rectangle (6.92,-6.92);
  \node[anchor=south, font=\footnotesize] at (3.46,0.25) {columns $j$};
  \node[anchor=east, rotate=90, font=\footnotesize] at (-0.3,-3.46) {rows $i$};
  % legend
  \fill[\nz] (8.0,-1.0) rectangle ++(0.7,0.7);
  \node[anchor=west] at (8.85,-0.65) {nonzero $\mat{K}_{ij}$};
  \draw[simGray!40] (8.0,-2.2) rectangle ++(0.7,0.7);
  \node[anchor=west] at (8.85,-1.85) {zero};
\end{tikzpicture}
\caption[Tridiagonal sparsity pattern]{The sparsity pattern of the assembled
stiffness matrix \cref{eq:tridiag} for $N=7$ interior nodes. Filled cells are the
nonzero entries; they sit only on the diagonal and the two adjacent off-diagonals,
because (\cref{fig:sparsity-support}) a hat couples only to its two
nearest-neighbour hats. Of the $49$ possible entries only $19$ are nonzero, and
the fraction shrinks as $N$ grows: a matrix this empty is never stored as a full
array.}
\label{fig:sparsity-grid}
\end{figure}

\subsection{Solving a small system}

Take $n=4$ cells, $h=\tfrac14$, so $N=3$ interior unknowns, and a constant load
$f(x)=1$. The load vector entries are
$\vb_i=\int_0^1 1\cdot\phi_i\dx=\text{area of the hat}=h=\tfrac14$ (each hat is a
triangle of base $2h$ and height $1$, area $h$). The system $\mat{K}\vx=\vb$ is,
multiplying \cref{eq:tridiag} by $h=\tfrac14$ to clear the prefactor,
\[
\frac1{1/4}
\begin{pmatrix} 2&-1&0\\-1&2&-1\\0&-1&2\end{pmatrix}
\begin{pmatrix}x_1\\x_2\\x_3\end{pmatrix}
=\begin{pmatrix}1/4\\1/4\\1/4\end{pmatrix}
\;\Longrightarrow\;
\begin{pmatrix} 2&-1&0\\-1&2&-1\\0&-1&2\end{pmatrix}
\begin{pmatrix}x_1\\x_2\\x_3\end{pmatrix}
=\begin{pmatrix}1/16\\1/16\\1/16\end{pmatrix}.
\]
Solving the tridiagonal system (eliminate downward, then back-substitute) gives
\[
x_1=\tfrac{3}{32}=0.09375,\quad
x_2=\tfrac{4}{32}=0.125,\quad
x_3=\tfrac{3}{32}=0.09375.
\]
The exact solution of $-u''=1$, $u(0)=u(1)=0$ is the parabola
$u(x)=\tfrac12 x(1-x)$, whose values at the nodes $x=\tfrac14,\tfrac12,\tfrac34$
are $\tfrac{3}{32},\tfrac14\!\cdot\!\tfrac34\!\cdot\!\tfrac12=\tfrac{3}{32}$\dots
let us check: $u(\tfrac14)=\tfrac12\cdot\tfrac14\cdot\tfrac34=\tfrac{3}{32}$,
$u(\tfrac12)=\tfrac12\cdot\tfrac12\cdot\tfrac12=\tfrac18=\tfrac{4}{32}$,
$u(\tfrac34)=\tfrac{3}{32}$. The finite-element coefficients match the exact
nodal values \emph{exactly}. This is no coincidence: for the 1-D Poisson problem
with the hat basis, the Galerkin solution is \emph{nodally exact}---a celebrated
special feature of this problem that does not survive into two dimensions or to
variable coefficients, but is a satisfying confirmation that the machinery works.

\begin{workedexample}{Is this $\mat{K}$ symmetric positive definite?}{spd-check}
We verify directly that the $3\times3$ stiffness above is SPD, the property
\cref{eq:spd} promised.

\emph{Symmetric.} By inspection $\mat{K}=\mat{K}^{\mathsf T}$: the
$(1,2)$ and $(2,1)$ entries are both $-1$, and so on. (We knew this must hold:
$a(u,v)=\int u'v'$ is symmetric in $u,v$.)

\emph{Positive definite.} Use the chain \cref{eq:spd}: for any
$\vy=(y_1,y_2,y_3)\neq\mathbf 0$,
\[
\vy^{\mathsf T}\mat{K}\vy
=2y_1^2+2y_2^2+2y_3^2-2y_1y_2-2y_2y_3
=y_1^2+(y_1-y_2)^2+(y_2-y_3)^2+y_3^2,
\]
a sum of squares that is zero only if $y_1=0$, $y_1=y_2$, $y_2=y_3$, $y_3=0$, i.e.
only if $\vy=\mathbf 0$. So $\vy^{\mathsf T}\mat{K}\vy>0$ for every nonzero $\vy$:
$\mat{K}$ is positive definite. (Equivalently, its leading principal minors
$2,\,3,\,4$ are all positive---Sylvester's criterion.) Because $\mat{K}$ is SPD,
$\mat{K}\vx=\vb$ has a unique solution and conjugate gradients (\cref{ch:cg}) is
guaranteed to find it.
\end{workedexample}

\section{Worked example: the mass matrix}

\begin{workedexample}{The 1-D hat mass matrix}{mass}
\label{wex:mass}
For the \emph{same} hat basis on the same uniform mesh, assemble the mass matrix
$\mat{M}_{ij}=\int_0^1\phi_i\phi_j\dx$---no derivatives this time. We need the
overlap integrals of the hats themselves.

\emph{Diagonal.} On its two support cells $\phi_i$ is a triangle of height $1$;
$\int\phi_i^2$ over one cell of length $h$ is, for a linear function running
$0\to1$, equal to $h/3$. Two cells give
\[
\mat{M}_{ii}=\int_{x_{i-1}}^{x_{i+1}}\phi_i^2\dx=\frac h3+\frac h3=\frac{2h}{3}.
\]

\emph{Off-diagonal.} On the shared cell $[x_i,x_{i+1}]$, $\phi_i$ runs $1\to0$
while $\phi_{i+1}$ runs $0\to1$; their product integrates to $h/6$:
\[
\mat{M}_{i,i+1}=\int_{x_i}^{x_{i+1}}\phi_i\,\phi_{i+1}\dx=\frac h6.
\]
Entries with $|i-j|\ge2$ vanish (disjoint support). Hence
\[
\mat{M}=\frac h6
\begin{pmatrix}
4 & 1 & & \\
1 & 4 & 1 & \\
 & 1 & 4 & \ddots\\
 & & \ddots & \ddots
\end{pmatrix}
=\frac h6\,\operatorname{tridiag}(1,\,4,\,1).
\]

\emph{Contrast with $\mat{K}$.} Both matrices are symmetric, tridiagonal, and
share the \emph{same} sparsity pattern (\cref{fig:sparsity-grid})---unsurprising,
since the nonzero pattern is fixed by which supports overlap, and both forms use
the same hats. But the \emph{signs} differ sharply: $\mat{K}$ has \emph{negative}
off-diagonals (a difference operator, $\diffop\phi_i\cdot\diffop\phi_j<0$ where
the slopes oppose), while $\mat{M}$ has \emph{positive} off-diagonals (an overlap
operator, $\phi_i\phi_j>0$ wherever the hats coexist). $\mat{M}$ is also SPD---it
is a Gram matrix of linearly independent functions---but it is much
better-conditioned than $\mat{K}$, which is why it appears on the ``easy'' side of
the generalized eigenproblem $\mat{K}\vx=\lambda\mat{M}\vx$
(\cref{ch:eigenmodes}).
\end{workedexample}

\section{The bilinear-form family is an operator family}
\label{sec:family}

We have so far assembled \emph{one} matrix from \emph{one} bilinear form. Real
analyses assemble several terms at once. The electrostatic stiffness is a single
diffusion term $(\varepsilon\Grad u,\Grad v)$; but the driven Maxwell operator
$\mat{A}(\omega)=\mat{K}+i\omega\mat{C}-\omega^2\mat{M}$ is a \emph{sum} of three,
and a lossy or multi-material problem adds still more. The structure that makes
this clean is the subject of \cref{ch:assembly}, but the essential idea belongs
here, because it is a direct consequence of the linearity of the form.

\subsection{Each term becomes one operator}

Every member of the bilinear-form family has the shape
$a_t(u,v)=(Q_t\,\diffop_t u,\diffop_t v)$ for some weight $Q_t$ and differential
operator $\diffop_t$---a permittivity-weighted gradient term, a
reluctivity-weighted curl term, a mass term, and so on. Sampling \emph{one} such
term on basis pairs produces \emph{one} matrix, exactly as in \cref{eq:KM}; call
it $\mat{A}(\text{term}_t)$. The full operator an analysis needs is the matrix of
its total form $a=\sum_t a_t$. Because sampling is linear in the form,
\begin{equation}
\mat{K}_{ij}=\Bigl(\sum_t a_t\Bigr)(\phi_j,\phi_i)
=\sum_t a_t(\phi_j,\phi_i)
=\sum_t \mat{A}(\text{term}_t)_{ij},
\qquad\text{i.e.}\qquad
\mat{K}=\sum_t \mat{A}(\text{term}_t).
\label{eq:assembly-sum}
\end{equation}
The global operator is a \emph{sum of per-term operators}. Adding a physical
effect---a boundary loss, a second material, the $i\omega\mat{C}$ damping---is
adding one more matrix to the sum, never rebuilding from scratch.

\begin{figure}[t]
\centering
\begin{tikzpicture}[node distance=6mm and 9mm, every node/.style={font=\footnotesize}]
  \node[opbox] (t1) {$\text{term}_1$\\$(\varepsilon\Grad u,\Grad v)$};
  \node[opbox, below=of t1] (t2) {$\text{term}_2$\\$(\sigma u, v)$};
  \node[opbox, below=of t2] (t3) {$\text{term}_3$\\$(\nu\Curl u,\Curl v)$};
  % per-term operators
  \node[stage, right=14mm of t1] (a1) {$\mat{A}_1$};
  \node[stage, right=14mm of t2] (a2) {$\mat{A}_2$};
  \node[stage, right=14mm of t3] (a3) {$\mat{A}_3$};
  \draw[flow] (t1) -- node[above,font=\scriptsize]{sample} (a1);
  \draw[flow] (t2) -- node[above,font=\scriptsize]{sample} (a2);
  \draw[flow] (t3) -- node[above,font=\scriptsize]{sample} (a3);
  % the sum node
  \node[product, right=20mm of a2] (K) {$\mat{K}=\sum_t\mat{A}_t$};
  \draw[flow] (a1) -- (K);
  \draw[flow] (a2) -- (K);
  \draw[flow] (a3) -- (K);
  % bracket label
  \node[font=\scriptsize\itshape, simGray, above=1mm of a1] {one operator per term};
\end{tikzpicture}
\caption[Assembly as a sum of term-operators]{The bilinear-form family becomes an
operator family. Each weak-form term $a_t(u,v)=(Q_t\diffop_t u,\diffop_t v)$
samples, on basis pairs, into its own matrix $\mat{A}_t$; the global operator of
the analysis is their sum, $\mat{K}=\sum_t\mat{A}_t$ (\cref{eq:assembly-sum}).
Because addition of operators is commutative and associative, the order in which
the terms are accumulated is immaterial---the structure of a \emph{commutative-
monoid fold} over the list of terms, which \cref{ch:assembly} makes precise.}
\label{fig:assembly-sum}
\end{figure}

\subsection{Assembly is a fold: a preview}

\Cref{eq:assembly-sum} is more than a formula; it is a \emph{computational
pattern}. We have a list of terms $[\text{term}_1,\dots,\text{term}_m]$, a way to
turn each into a matrix, and a way to combine matrices (add them). Combining a
list of things with an associative, commutative operation that has an identity
(the zero matrix) is the canonical setting for a \emph{fold}---the higher-order
function that walks a list accumulating a result. In the pseudo-language of
Part~III we will write assembly as
\begin{equation}
\mat{K}=\texttt{fe\_assemble}(\textit{space},\textit{terms})
=\textstyle\sum_{t\in\textit{terms}}\mat{A}(\textit{space},t),
\label{eq:fe-assemble}
\end{equation}
literally a fold of the per-term assembly map over the immutable list of terms.
Two features of this fold are worth flagging now, because they are properties of
the \emph{form}, not of any particular code:

\begin{itemize}[nosep]
\item \textbf{Order independence.} Operator addition is commutative and
associative, so $\sum_t\mat{A}_t$ does not depend on the order the terms are
listed. The fold is over a \emph{commutative monoid}; the term list is really an
unordered bag.
\item \textbf{The list-homomorphism intuition.} Assembling a concatenated term
list equals assembling the pieces and adding:
$\texttt{fe\_assemble}(\textit{space},\,T_1\!+\!\!+\,T_2)=
\texttt{fe\_assemble}(\textit{space},T_1)+\texttt{fe\_assemble}(\textit{space},T_2)$.
A function that turns list-concatenation into a combine on results is a
\emph{list homomorphism}, the structure that makes assembly splittable, mergeable,
and---on a parallel machine---divisible across processors.
\end{itemize}

This is why \cref{ch:assembly} can treat assembly as one combinator rather than a
tangle of nested loops, and why the matrix-free operator of \cref{ch:matrixfree}
can apply the global $\mat{K}$ without ever forming it: applying a sum of
operators is applying each and adding. The bilinear-form family, sampled, is an
operator family, folded.

\section{How good is the Galerkin solution?}
\label{sec:cea}

We have a recipe that produces \emph{some} $u_h$ in $V_h$. Is it any good? The
reassuring answer, made precise by a one-line argument of Jean Céa, is that the
Galerkin solution is \emph{quasi-optimal}: it is, up to a fixed constant, as
accurate as the very best approximation the subspace $V_h$ could possibly offer.

\subsection{Galerkin orthogonality}

Start from a small but powerful fact. The exact solution $u$ satisfies
$a(u,v)=\ell(v)$ for all $v\in V$, and in particular for all $v_h\in V_h\subset V$.
The Galerkin solution $u_h$ satisfies $a(u_h,v_h)=\ell(v_h)$ for all $v_h\in V_h$.
Subtracting,
\begin{equation}
a(u-u_h,\,v_h)=0
\qquad\text{for all } v_h\in V_h.
\label{eq:galerkin-orthog}
\end{equation}
The error $u-u_h$ is invisible to every test function in $V_h$: in the geometry
the form $a$ defines, the error is \emph{orthogonal} to the whole subspace. This
is the precise content of \cref{fig:galerkin-projection}---$u_h$ is the orthogonal
projection of $u$ onto $V_h$ in the energy inner product---and it is the engine of
the next result.

\subsection{Céa's lemma}

\begin{lemma}[Céa, quasi-optimality]
Suppose $a$ is bounded, $\abs{a(u,v)}\le M\norm{u}\,\norm{v}$, and coercive,
$a(v,v)\ge\alpha\norm{v}^2$. Then the Galerkin error is controlled by the best
approximation error in $V_h$:
\begin{equation}
\norm{u-u_h}\le\frac{M}{\alpha}\,\min_{v_h\in V_h}\norm{u-v_h}.
\label{eq:cea}
\end{equation}
\end{lemma}

\begin{proof}
Let $v_h\in V_h$ be \emph{any} element of the subspace. Using coercivity, then
Galerkin orthogonality \cref{eq:galerkin-orthog} to subtract the cross term
$a(u-u_h,v_h-u_h)=0$ (legal because $v_h-u_h\in V_h$), then boundedness:
\[
\alpha\norm{u-u_h}^2
\le a(u-u_h,\,u-u_h)
= a(u-u_h,\,u-v_h)
\le M\,\norm{u-u_h}\,\norm{u-v_h}.
\]
Divide both sides by $\alpha\norm{u-u_h}$ (zero error is trivially optimal) to get
$\norm{u-u_h}\le(M/\alpha)\norm{u-v_h}$. Since $v_h$ was arbitrary, take the
minimum over $v_h\in V_h$.
\end{proof}

Read \cref{eq:cea} carefully, because it is the deepest single fact in this
chapter. The right-hand side, $\min_{v_h}\norm{u-v_h}$, is a question of
\emph{approximation theory} alone: how well can the functions in $V_h$ approximate
$u$, ignoring the differential equation entirely? Céa's lemma says the Galerkin
method---which knows nothing of $u$ and only solves $\mat{K}\vx=\vb$---comes
within a constant factor $M/\alpha$ of that best-possible answer
(\cref{fig:cea}). The method does not need to be clever about $u$; it only needs a
subspace $V_h$ that \emph{can} approximate $u$ well. The job of choosing such a
subspace is the job of the basis, which returns us to the chapter's refrain.

\begin{figure}[t]
\centering
\begin{tikzpicture}[scale=1.0, every node/.style={font=\footnotesize}]
  % subspace line V_h
  \draw[draw=simGreen, fill=simBgGreen, rounded corners=4pt]
    (-3.2,-0.35) rectangle (3.2,0.35);
  \draw[simGreen!60!simInk, thick] (-3.0,0)--(3.0,0);
  \node[simGreen!70!simInk, left] at (-3.05,0) {$V_h$};
  % exact solution above
  \fill[simRust] (0.0,1.9) circle (1.8pt);
  \node[simRust, above] at (0.0,2.0) {$u$};
  % best approx (foot of perpendicular)
  \fill[simInk] (0.0,0) circle (1.8pt);
  \node[simInk, below right=-1pt] at (0.05,-0.05) {$v_h^\star$ (best)};
  % galerkin solution nearby
  \fill[simBlue] (1.3,0) circle (1.8pt);
  \node[simBlue, below] at (1.3,-0.12) {$u_h$ (Galerkin)};
  % best distance (perpendicular)
  \draw[flow, simInk] (0.0,1.9) -- (0.0,0.06);
  \node[simInk, right, font=\scriptsize] at (0.06,1.0) {$\min_{v_h}\|u-v_h\|$};
  % galerkin distance
  \draw[flow, simBlue] (0.0,1.9) -- (1.3,0.06);
  \node[simBlue, right=1pt, font=\scriptsize] at (0.78,1.05) {$\|u-u_h\|$};
  % right-angle marker for best approx
  \draw[simInk] (0.0,0.18)--(0.18,0.18)--(0.18,0);
\end{tikzpicture}
\caption[Céa's quasi-optimality]{Céa's lemma, geometrically. The best the
subspace $V_h$ can do is the foot of the perpendicular $v_h^\star$, at distance
$\min_{v_h}\norm{u-v_h}$ from the exact solution $u$. The Galerkin solution $u_h$
need not be that closest point (it is closest in the \emph{energy} norm, not the
norm drawn), but \cref{eq:cea} guarantees it is no worse than a fixed multiple
$M/\alpha$ of the best distance. The Galerkin method is therefore only as good---
and only as bad---as the approximating power of the basis.}
\label{fig:cea}
\end{figure}

\subsection{Convergence: \texorpdfstring{$h$}{h}- and \texorpdfstring{$p$}{p}-refinement}

Quasi-optimality turns convergence questions into approximation questions, and
approximation theory then supplies the rates. Two knobs control how well $V_h$
can approximate, and refining either drives the error down.

\emph{$h$-refinement} shrinks the mesh spacing $h$, adding more, smaller elements.
For a basis of piecewise polynomials of degree $p$ on a mesh of size $h$,
approximating a smooth $u$, standard interpolation estimates give
$\min_{v_h}\norm{u-v_h}=\mathcal{O}(h^{p})$ in the energy norm, so by Céa
\begin{equation}
\norm{u-u_h}=\mathcal{O}(h^{p}).
\label{eq:hrate}
\end{equation}
With the linear hats ($p=1$) of our worked example, halving $h$ halves the energy
error; the error in the solution value itself converges one order faster,
$\mathcal{O}(h^{p+1})$. More elements, more accuracy, at a predictable rate.

\emph{$p$-refinement} keeps the mesh but raises the polynomial degree $p$ of the
basis on each element. For smooth $u$ this is dramatically more effective: the
error falls \emph{faster than any fixed power of $1/p$}---so-called spectral or
exponential convergence---because high-degree polynomials approximate smooth
functions extraordinarily well. Palace exploits both knobs, and combines them in
the $p$-multigrid hierarchy of \cref{ch:multigrid}. Which knob to turn is a
judgement about the smoothness of the solution: $h$-refinement is robust near
singularities (sharp corners, material interfaces), $p$-refinement wins where the
field is smooth.

\begin{remark}
The constant $M/\alpha$ in \cref{eq:cea} is the form's \emph{condition number} in
disguise---the ratio of its boundedness constant to its coercivity constant. The
same ratio reappears as the condition number $\kappa(\mat{K})$ of the assembled
matrix, which (\cref{ch:cg}) sets the convergence rate of conjugate gradients.
A form that is barely coercive ($\alpha$ small) gives both a loose error bound and
a slow solver---another instance of a continuous property dictating discrete
behaviour.
\end{remark}

\section{The basis is the whole game}
\label{sec:basis-is-everything}

Step back and notice how little of Galerkin's method depended on the
\emph{physics}. The recipe $\mat{K}_{ij}=a(\phi_j,\phi_i)$, $\vb_i=\ell(\phi_i)$,
solve $\mat{K}\vx=\vb$, is identical for the Poisson equation, the curl--curl
equation, and every analysis in this book. What changes from one problem to the
next, beyond the form $a$, is the \emph{basis} $\{\phi_i\}$---and the basis alone
determines:

\begin{itemize}[nosep]
\item the \textbf{size} $N$ of the system (one unknown per basis function);
\item the \textbf{sparsity} of $\mat{K}$ (which supports overlap,
\cref{sec:sparsity});
\item the \textbf{accuracy}, through the best-approximation term in Céa's lemma
(\cref{sec:cea});
\item and, crucially, \textbf{which physical fields the method can represent
faithfully.}
\end{itemize}

That last point is where the rest of Part~II is heading. A scalar potential $V$
is well represented by the continuous, nodal hat functions of this chapter---an
$\Hone$-conforming basis, in the language of \cref{ch:functionspaces}. But an
electric field $\vE$ is a \emph{vector} field with a different continuity
requirement (its tangential component, not its full value, is continuous across
material interfaces), and forcing nodal hats onto it produces spurious,
unphysical modes. The fix is a different basis entirely---the edge elements of
$\Hcurl$, whose degrees of freedom live on mesh \emph{edges} rather than nodes.
Which basis a field needs is dictated by the differential operator $\diffop$ in
its form: $\Grad$ wants nodes, $\Curl$ wants edges, $\Div$ wants faces. Building
those bases, and the reference elements they are constructed on, is the work of
\cref{ch:functionspaces,ch:refelement}. Galerkin's method is the same for all of
them; only the menu of functions changes.

\summarysection
Galerkin's method makes the weak form computable by replacing the
infinite-dimensional trial/test space $V$ with a finite-dimensional subspace
$V_h=\operatorname{span}\{\phi_1,\dots,\phi_N\}$ and requiring the weak equation
only against the $N$ basis functions. Expanding the unknown $u_h=\sum_j u_j\phi_j$
turns this into a square linear system $\mat{K}\vx=\vb$ with
$\mat{K}_{ij}=a(\phi_j,\phi_i)$ and $\vb_i=\ell(\phi_i)$---the matrix \emph{is}
the bilinear form sampled on basis pairs, and the mass matrix
$\mat{M}_{ij}=\int\phi_i\phi_j$ is the same construction with the identity
operator. The matrix inherits its properties from the form: symmetry (enabling
conjugate gradients), positive-definiteness from coercivity (guaranteeing a unique
solution), and---from the \emph{local support} of the basis---sparsity, which is
why we solve by iteration. The bilinear-form family is an operator family: each
term $a_t=(Q_t\diffop_t u,\diffop_t v)$ assembles into one matrix and the global
operator is their sum $\mat{K}=\sum_t\mat{A}_t$, a commutative-monoid fold and
list homomorphism over the term list---the assembly-as-a-fold of \cref{ch:assembly}.
Céa's lemma shows the Galerkin solution is quasi-optimal: within a constant
$M/\alpha$ of the best approximation $V_h$ can provide, so convergence ($h$- and
$p$-refinement) is governed by the approximating power of the basis. The basis,
in the end, is the whole game.

\furtherreading
The variational and approximation theory of this chapter---Galerkin orthogonality,
Céa's lemma, the $h$- and $p$-estimates, and the SPD structure of stiffness
matrices---is developed rigorously in Brenner and Scott~\citep{brennerscott2008},
the standard graduate reference for the mathematics of finite elements; their
treatment of the bilinear-form framework and approximation rates expands every
result we stated. For the linear-algebra side---why sparsity and conditioning
decide between conjugate gradients and GMRES, and how the matrix's spectral
properties govern iterative convergence---Trefethen and Bau~\citep{trefethenbau1997}
is the classic and exceptionally readable account. For the engineering
practice of assembling and solving finite-element systems for electromagnetics
specifically, including vector ($\Hcurl$) bases and the families we turn to next,
Jin~\citep{jin2014} is the comprehensive reference.

\exercisessection
\begin{enumerate}[nosep]
  \item \textbf{One more node.} Redo the worked example of
  \cref{sec:worked-stiffness} with $n=5$ cells ($h=\tfrac15$, $N=4$ unknowns) and
  constant load $f=1$. Write down the $4\times4$ system $\mat{K}\vx=\vb$ and solve
  it. Confirm that the coefficients again equal the exact nodal values of
  $u=\tfrac12 x(1-x)$.

  \item \textbf{Variable load.} Keep $n=4$, $h=\tfrac14$, but take
  $f(x)=x$. Compute the load vector $\vb_i=\int_0^1 x\,\phi_i(x)\dx$ (the
  integral of $x$ against each hat) and solve $\mat{K}\vx=\vb$ with the same
  $\mat{K}$ as in the text. Why did $\mat{K}$ not change when the right-hand side
  did?

  \item \textbf{Counting nonzeros.} For the 1-D hat basis on $N$ interior nodes,
  how many nonzero entries does $\mat{K}$ have? Express it as a function of $N$ and
  confirm it is $\mathcal{O}(N)$, not $\mathcal{O}(N^2)$. How many nonzeros would
  a 2-D mesh in which each node couples to six neighbours have per row?

  \item \textbf{Symmetry from the form.} Show that \emph{any} form of the family
  $a(u,v)=\int_\dom Q\,\diffop u\cdot\diffop v$ with scalar $Q$ assembles a
  symmetric matrix, directly from $\mat{K}_{ij}=a(\phi_j,\phi_i)$. Then give an
  example of a form (a first-derivative ``advection'' term $\int u'\,v$) whose
  matrix is \emph{not} symmetric, and identify which step of the symmetry argument
  fails.

  \item \textbf{Mass-matrix positive-definiteness.} Using the chain
  \cref{eq:spd}, argue that the mass matrix $\mat{M}_{ij}=\int\phi_i\phi_j$ is
  positive definite whenever the basis functions $\{\phi_i\}$ are linearly
  independent. (Hint: $\vy^{\mathsf T}\mat{M}\vy=\int(\sum_j y_j\phi_j)^2$. When is
  this zero?)

  \item \textbf{Best approximation versus Galerkin.} For the $n=4$ Poisson
  example, the Galerkin solution was \emph{nodally exact}. The best $L^2$
  approximation of $u=\tfrac12 x(1-x)$ in $V_h$ minimizes $\int(u-v_h)^2$ and is
  generally \emph{not} nodally exact. Without computing it, explain using Céa's
  lemma and Galerkin orthogonality why the two need not coincide, and which inner
  product the Galerkin solution \emph{is} optimal in.

  \item \textbf{Assembly as a fold.} An analysis has three weak-form terms with
  per-term matrices $\mat{A}_1,\mat{A}_2,\mat{A}_3$. Write the global operator
  three different ways---$\mat{A}_1+(\mat{A}_2+\mat{A}_3)$,
  $(\mat{A}_1+\mat{A}_2)+\mat{A}_3$, and $\mat{A}_2+\mat{A}_3+\mat{A}_1$---and
  argue they are equal. Which algebraic properties of operator addition did you
  use, and why do they make assembly a \emph{commutative-monoid} fold rather than
  merely a list fold?

  \item \textbf{Why not a global basis?} Suppose you chose the global sine basis
  $\phi_k(x)=\sin(k\pi x)$ on $(0,1)$ instead of hats. (a)~Show every pair of
  supports overlaps, so the sparsity argument of \cref{sec:sparsity} fails.
  (b)~Compute $\mat{K}_{ij}=\int_0^1\phi_i'\phi_j'$ for this basis and show it
  comes out \emph{diagonal} anyway. (c)~Explain why this pleasant accident does
  not save the finite-element method in general, and what is lost by giving up
  local support.
\end{enumerate}
